\chapter{Motion Control}
\label{ch:motion-control}

Following a desired trajectory, or reaching a desired configuration,
cannot be done by simply applying the inputs computed from the flat
outputs: disturbances and modeling errors would accumulate without
bound. A feedback action is needed. In previous robotics courses, the
dynamic model of manipulators was controlled through joint torques.
For wheeled mobile robots we can work at the kinematic level, because
the dynamic model is simpler, can be canceled via feedback, and the
wheels are equipped with low-level velocity loops that are reliable.

\begin{figure}[ht]
    \centering
    \includegraphics[width=0.8\linewidth]{images/motion_control_scheme.png}
    \caption{Actual control scheme.}
    \label{fig:control-scheme-actual}
\end{figure}

\begin{figure}[ht]
    \centering
    \includegraphics[width=0.75\linewidth]{images/motion_control_scheme2.png}
    \caption{Simplified control scheme used for design.}
    \label{fig:control-scheme-design}
\end{figure}

Two control tasks are considered in this course:
\begin{itemize}
    \item \textbf{Trajectory tracking}: the robot must converge to a
          desired Cartesian trajectory \( (x_d(t), y_d(t)) \). Treated
          in this chapter.
    \item \textbf{Posture regulation}: the robot must converge to a
          desired fixed configuration \( q_d \). Treated in
          \cref{ch:regulation}.
\end{itemize}
Trajectory tracking is more relevant in practice, because moving along
planned trajectories reduces the risk of collisions. Posture
regulation, on the other hand, does not require any planning.

\section{Trajectory Tracking for the Unicycle}
\label{sec:trajectory-tracking}

Let \( (x_d(t), y_d(t)) \) be an admissible desired trajectory for the
unicycle. Admissibility means that there exist inputs \( v_d, \omega_d \)
such that
\begin{equation}
\label{eq:desired-unicycle}
    \begin{cases}
        \dot{x}_d = v_d \cos\theta_d, \\
        \dot{y}_d = v_d \sin\theta_d, \\
        \dot{\theta}_d = \omega_d.
    \end{cases}
\end{equation}
Thanks to the flatness property (\cref{sec:differential-flatness}),
from the desired trajectory \( (x_d(t), y_d(t)) \) we can reconstruct
the full desired state and the required inputs:
\begin{align}
    \theta_d(t) &= \operatorname{atan2}(\dot{y}_d(t), \dot{x}_d(t)) + k\pi,
                   \quad k \in \{0, 1\}, \\
    v_d(t) &= \pm\sqrt{\dot{x}_d^2(t) + \dot{y}_d^2(t)}, \\
    \omega_d(t) &= \frac{\ddot{y}_d(t) \dot{x}_d(t)
                        - \ddot{x}_d(t) \dot{y}_d(t)}
                       {\dot{x}_d^2(t) + \dot{y}_d^2(t)}.
\end{align}
An admissible trajectory requires \( (\dot{x}_d(t), \dot{y}_d(t)) \) to
be continuous over time. The desired state trajectory \( q_d(t) \)
computed from the flat outputs is used to define the state error, and
the resulting block diagram is in \cref{fig:control-scheme-flatness}.

\begin{figure}[ht]
    \centering
    \includegraphics[width=0.8\linewidth]{images/motion_control_scheme3.png}
    \caption{Feedforward from flatness, feedback from state error.}
    \label{fig:control-scheme-flatness}
\end{figure}

\subsection{Error Dynamics}
\label{subsec:error-dynamics}

The objective is to drive the state error to zero. The error is
expressed in a frame rotated by \( \theta \), with the \( x \)-axis
parallel to the sagittal axis of the robot:
\begin{equation}
\label{eq:error-rotated}
    e = \begin{pmatrix} e_1 \\ e_2 \\ e_3 \end{pmatrix}
    =
    \begin{pmatrix}
        \cos\theta & \sin\theta & 0 \\
        -\sin\theta & \cos\theta & 0 \\
        0 & 0 & 1
    \end{pmatrix}
    \begin{pmatrix}
        x_d - x \\ y_d - y \\ \theta_d - \theta
    \end{pmatrix}.
\end{equation}
The pair \( (e_1, e_2) \) is the Cartesian error
\( e_p = (x_d - x,\, y_d - y)^\top \) expressed in the rotated frame,
as in \cref{fig:error-rotated}. Differentiating \( e \) and using the
unicycle kinematic model together with the desired one, we obtain the
\textbf{error dynamics}:
\begin{align}
\label{eq:error-dynamics}
    \dot{e}_1 &= v_d \cos e_3 - v + e_2 \omega, \\
    \dot{e}_2 &= v_d \sin e_3 - e_1 \omega, \\
    \dot{e}_3 &= \omega_d - \omega.
\end{align}
The error dynamics is \textbf{nonlinear} and \textbf{time-varying},
since it depends on the reference inputs \( v_d(t), \omega_d(t) \).

\begin{figure}[ht]
    \centering
    \includegraphics[width=0.45\linewidth]{images/error.eps}
    \caption{Cartesian error expressed in the frame rotated by
    \( \theta \).}
    \label{fig:error-rotated}
\end{figure}

\subsection{Approximate Linearization}
\label{subsec:approximate-linearization}

Before applying it to the unicycle, recall the idea of stabilization
via approximate linearization: to stabilize a nonlinear system at an
equilibrium, we stabilize its approximate linearization around it.
Consider a generic nonlinear system
\[
\dot{x} = \eta(x, u) = f(x) + G(x) u,
\qquad x \in \mathbb{R}^n,\ u \in \mathbb{R}^m,
\]
and assume, without loss of generality, that the origin is the desired
\textbf{unforced equilibrium}, i.e.\ \( \eta(0, 0) = 0 \), or
equivalently \( f(0) = 0 \). The first-order Taylor expansion of
\( \eta \) around \( x = 0, u = 0 \) is
\[
\eta(x, u) \approx \underbrace{\eta(0, 0)}_{=\,0}
    + \left. \frac{\partial \eta}{\partial x} \right|_{0}
      x
    + \left. \frac{\partial \eta}{\partial u} \right|_{0}
      u,
\]
so the \textbf{approximate linearization} of the system at the origin is
\[
\dot{x} = A x + B u,
\]
where \( A \) and \( B \) are the two Jacobians evaluated at the
origin. With the linear feedback \( u = K x \),
\[
\dot{x} = (A + BK) x.
\]
If \( K \) is chosen so that \( A + BK \) is \textbf{Hurwitz}
(certainly possible if \( (A, B) \) is controllable), the approximate
linearization is asymptotically stable.

\begin{theorem}[Lyapunov Indirect Method]
\label{thm:lyapunov-indirect}
If the approximate linearization of a nonlinear system at an
equilibrium is asymptotically stable, then the equilibrium is
\textbf{locally} asymptotically stable for the nonlinear system.
\end{theorem}

\paragraph{Application to the unicycle.}
In the form \cref{eq:error-dynamics}, the origin is \textbf{not} an
unforced equilibrium: for \( e = 0 \), the terms \( v_d - v \) and
\( \omega_d - \omega \) must vanish to have \( \dot{e} = 0 \). This is
fixed by the (invertible) input transformation
\begin{align}
    w_1 &= v_d \cos e_3 - v, \\
    w_2 &= \omega_d - \omega.
\end{align}
Substituting \( \omega = \omega_d - w_2 \), the error dynamics
\( \dot{e} = \eta(e, w) \) takes the form
\begin{equation}
    \dot{e} = \underbrace{
    \begin{pmatrix}
        \omega_d e_2 \\
        -\omega_d e_1 + v_d \sin e_3 \\
        0
    \end{pmatrix}}_{f(e)}
    +
    \underbrace{
    \begin{pmatrix}
        1 & -e_2 \\
        0 & e_1 \\
        0 & 1
    \end{pmatrix}}_{G(e)}
    \begin{pmatrix} w_1 \\ w_2 \end{pmatrix},
\end{equation}
which is at rest for \( e = 0, w = 0 \), i.e.\
\( \eta(0, 0) = 0 \). Both the drift and the input term are
nonlinear, and the drift is in general time-varying. The Jacobians are
\[
\frac{\partial \eta}{\partial e} =
\begin{pmatrix}
    0 & \omega_d - w_2 & 0 \\
    -\omega_d + w_2 & 0 & v_d \cos e_3 \\
    0 & 0 & 0
\end{pmatrix},
\qquad
\frac{\partial \eta}{\partial w} =
\begin{pmatrix}
    1 & -e_2 \\
    0 & e_1 \\
    0 & 1
\end{pmatrix},
\]
and evaluating at \( e = 0, w = 0 \):
\begin{equation}
\label{eq:A-tracking}
    A(t) =
    \begin{pmatrix}
        0 & \omega_d & 0 \\
        -\omega_d & 0 & v_d \\
        0 & 0 & 0
    \end{pmatrix},
    \qquad
    B =
    \begin{pmatrix}
        1 & 0 \\
        0 & 0 \\
        0 & 1
    \end{pmatrix}.
\end{equation}
The pair \( (A(t), B) \) is controllable iff the matrix
\[
\begin{pmatrix} B & A B \end{pmatrix}
\]
has full rank. In this case,
\[
\operatorname{rank}
\begin{pmatrix}
    1 & 0 & 0 & 0 \\
    0 & 0 & -\omega_d & v_d \\
    0 & 1 & 0 & 0
\end{pmatrix}
= 3,
\]
which holds if \( \omega_d \neq 0 \) or \( v_d \neq 0 \). The system
loses controllability when the unicycle stops.

Note that \( A \) is actually \( A(t) \), because the reference inputs
\( v_d, \omega_d \) depend on time: the linear approximation is still
a \textbf{time-varying} system. Define the linear feedback
\[
w = \begin{pmatrix} w_1 \\ w_2 \end{pmatrix}
  = \underbrace{
    \begin{pmatrix}
        -k_1 & 0 & 0 \\
        0 & -k_2 & -k_3
    \end{pmatrix}}_{K}
    \begin{pmatrix} e_1 \\ e_2 \\ e_3 \end{pmatrix}
  = K e.
\]
The closed-loop error dynamics is
\[
\dot{e} = (A(t) + BK) e = A_{cl}(t) e =
\begin{pmatrix}
    -k_1 & \omega_d & 0 \\
    -\omega_d & 0 & v_d \\
    0 & -k_2 & -k_3
\end{pmatrix} e,
\]
still time-varying. Setting
\[
k_1 = k_3 = 2 \zeta a, \qquad
k_2 = \frac{a^2 - \omega_d^2}{v_d},
\]
with \( a > 0 \) and \( \zeta \in (0, 1) \), the characteristic
polynomial of \( A_{cl}(t) \) becomes time-invariant and Hurwitz:
\[
p(\lambda) = (\lambda + 2 \zeta a)
             (\lambda^2 + 2 \zeta a \lambda + a^2).
\]
It has one real negative eigenvalue and two complex eigenvalues with
damping \( \zeta \) and natural frequency \( a \).

\begin{nb}[Caveats of Approximate Linearization]
Having the eigenvalues of \( A_{cl}(t) \) fixed at every instant does
\emph{not} guarantee asymptotic stability of a time-varying system,
unless \( v_d \) and \( \omega_d \) are constant (rectilinear and
circular trajectories). Even in that case, asymptotic stability is
only \emph{local}, because it comes from \cref{thm:lyapunov-indirect}.
Moreover, \( k_2 \to \infty \) as \( v_d \to 0 \): the controller can
only be used with \textbf{persistent} Cartesian trajectories (stops
are not allowed).
\end{nb}

The \textbf{actual} velocity inputs are recovered from the input
transformation:
\begin{align}
    v &= v_d \cos e_3 - w_1 = v_d \cos e_3 + k_1 e_1, \\
    \omega &= \omega_d - w_2 = \omega_d + k_2 e_2 + k_3 e_3.
\end{align}
As \( e \to 0 \), \( (v, \omega) \to (v_d, \omega_d) \): the control
reduces to the pure feedforward term. The controller is \textbf{static}
and based on the state error. It needs \( v_d, \omega_d \) (computed
via flatness) and \( \theta \), both for the error rotation and for
the input transformation.

\textbf{Global stability} can be guaranteed by a nonlinear version of
the feedback:
\begin{align}
    w_1 &= -k_1(v_d, \omega_d) e_1, \\
    w_2 &= -k_2 v_d \frac{\sin e_3}{e_3} e_2
           - k_3(v_d, \omega_d) e_3,
\end{align}
with \( k_2 > 0 \) constant and \( k_1, k_3 \) bounded, positive, with
bounded derivatives. The actual inputs are recovered from the same
input transformation.

\subsection{Output Error Feedback: Static}
\label{subsec:output-error-static}

A different approach builds the feedback from the \textbf{output
(Cartesian) error} only, without computing a desired state trajectory
\( q_d(t) \), using the velocity \( \dot{y}_d \) along the reference
trajectory as the \textbf{feedforward term}. Here
\( y = (x, y)^\top \) is the Cartesian position of the unicycle and
\( y_d \) the desired one. The controller is designed by
\textbf{exact input-output linearization}.

\paragraph{Exact input-output linearization.}
Consider a driftless nonlinear system with the same number of inputs
and outputs:
\begin{align}
    \dot{x} &= G(x) u, \qquad x \in \mathbb{R}^n,\ u \in \mathbb{R}^m,
    \\
    y &= h(x), \qquad y \in \mathbb{R}^m.
\end{align}
Differentiating the output,
\[
\dot{y} = \frac{\partial h}{\partial x} \dot{x}
        = \underbrace{\frac{\partial h}{\partial x} G(x)}_{T(x)} u
        = T(x) u,
\]
where \( T \) is an \( m \times m \) matrix called the
\textbf{decoupling matrix}. If \( T \) is invertible, the
\textbf{linearizing feedback}
\[
u = T^{-1}(x) w
\]
yields \( \dot{y} = w \): an \textbf{exactly linear} map from the new
inputs \( w \) to the derivative of the outputs. Given a reference
output \( y_d(t) \), the output error \( e = y_d - y \) evolves as
\[
\dot{e} = \dot{y}_d - \dot{y} = \dot{y}_d - w.
\]
With the choice \( w = \dot{y}_d + K e \) (feedforward plus
proportional feedback),
\[
\dot{e} = -K e,
\]
which is \textbf{globally exponentially stable} if the eigenvalues of
\( K \) lie in the right half-plane. The final control law is
\[
u = T^{-1}(x) \big( \dot{y}_d + K e \big).
\]

\paragraph{Singular decoupling matrix.}
For the unicycle with output \( y = (x, y)^\top \), the decoupling
matrix is singular:
\[
\begin{pmatrix} \dot{x} \\ \dot{y} \end{pmatrix}
=
\begin{pmatrix}
    \cos\theta & 0 \\
    \sin\theta & 0
\end{pmatrix}
\begin{pmatrix} v \\ \omega \end{pmatrix}.
\]
Hence exact input-output linearization is \textbf{not possible} for
the output \( (x, y) \).

\paragraph{Displaced output.}
The solution is to move the controlled output from the wheel contact
point to a point \( B \) along the sagittal axis, at distance \( b \):
\[
\begin{pmatrix} x_B \\ y_B \end{pmatrix}
=
\begin{pmatrix}
    x + b \cos\theta \\
    y + b \sin\theta
\end{pmatrix}.
\]
Differentiating,
\[
\begin{pmatrix} \dot{x}_B \\ \dot{y}_B \end{pmatrix}
=
\underbrace{
\begin{pmatrix}
    \cos\theta & -b \sin\theta \\
    \sin\theta & b \cos\theta
\end{pmatrix}}_{T(\theta),\ \det T = b}
\begin{pmatrix} v \\ \omega \end{pmatrix}.
\]
If \( b \neq 0 \), set
\[
\begin{pmatrix} v \\ \omega \end{pmatrix}
= T^{-1}(\theta)
  \begin{pmatrix} w_1 \\ w_2 \end{pmatrix}
=
\begin{pmatrix}
    \cos\theta & \sin\theta \\
    -\dfrac{\sin\theta}{b} & \dfrac{\cos\theta}{b}
\end{pmatrix}
\begin{pmatrix} w_1 \\ w_2 \end{pmatrix},
\]
obtaining the input-output linearized system
\begin{align}
    \dot{x}_B &= w_1, \\
    \dot{y}_B &= w_2, \\
    \dot{\theta} &= \frac{w_2 \cos\theta - w_1 \sin\theta}{b}.
\end{align}
\textbf{Global exponential convergence} of \( x_B, y_B \) to the
desired trajectory is achieved with
\begin{align}
    w_1 &= \dot{x}_d + k_x (x_d - x_B), \\
    w_2 &= \dot{y}_d + k_y (y_d - y_B),
\end{align}
with \( k_x, k_y > 0 \). The orientation \( \theta \) is \textbf{not}
controlled by this scheme: there is a \textbf{zero dynamics}, the last
equation above. Still, \( \theta \) must evolve as dictated by
flatness.

An advantage of this approach is that the desired trajectory for
\( B \) can be \textbf{arbitrary}: square corners may be included. The
controller is \textbf{static}, based on the output error, it needs
\( \dot{y}_d \), and it needs \( x, y, \theta \) for the output
reconstruction and \( \theta \) for the input transformation.

\subsection{Output Error Feedback: Dynamic}
\label{subsec:output-error-dynamic}

Rather than displacing the output to \( B \), one can keep the (flat)
output \( y = (x, y)^\top \) and achieve exact linearization with a
\textbf{dynamic compensator}. Perform a \textbf{dynamic extension} on
the driving velocity channel: treat \( v \) as a state and its
acceleration \( a \) as a new input,
\[
\dot{v} = a.
\]
Now differentiate the output further:
\begin{align}
    \ddot{x} &= a \cos\theta - v \sin\theta\,\omega, \\
    \ddot{y} &= a \sin\theta + v \cos\theta\,\omega.
\end{align}
The new decoupling matrix is
\[
\begin{pmatrix} \ddot{x} \\ \ddot{y} \end{pmatrix}
=
\underbrace{
\begin{pmatrix}
    \cos\theta & -v \sin\theta \\
    \sin\theta & v \cos\theta
\end{pmatrix}}_{T(v, \theta),\ \det T = v}
\begin{pmatrix} a \\ \omega \end{pmatrix},
\]
which is nonsingular if \( v \neq 0 \). Under this assumption,
\[
\begin{pmatrix} a \\ \omega \end{pmatrix}
= T^{-1}(v, \theta)
  \begin{pmatrix} w_1 \\ w_2 \end{pmatrix}
=
\begin{pmatrix}
    \cos\theta & \sin\theta \\
    -\dfrac{\sin\theta}{v} & \dfrac{\cos\theta}{v}
\end{pmatrix}
\begin{pmatrix} w_1 \\ w_2 \end{pmatrix},
\]
and the system becomes \( \ddot{x} = w_1 \),
\( \ddot{y} = w_2 \). \textbf{Global exponential convergence} of
\( x, y \) to the desired trajectory is achieved with
\begin{align}
    w_1 &= \ddot{x}_d + k_{p,x} (x_d - x) + k_{d,x} (\dot{x}_d - \dot{x}), \\
    w_2 &= \ddot{y}_d + k_{p,y} (y_d - y) + k_{d,y} (\dot{y}_d - \dot{y}),
\end{align}
with \( k_{p,x}, k_{p,y}, k_{d,x}, k_{d,y} > 0 \). Since the original
flat output is controlled, there is \textbf{no zero dynamics} with
this approach.

The desired trajectory must be twice differentiable and
\textbf{persistent}, i.e.\ \( v_d \neq 0 \) always (the singularity is
structural in nonholonomic systems). The controller is \textbf{dynamic}
(the velocity \( v \) is the state of the compensator, obtained by
integrating \( a \)), based on the output error. It needs
\( \dot{y}_d \) and \( \ddot{y}_d \), and it needs \( x, y \) for the
error computation and \( \theta \) for the input transformation.

\subsection{Comparison of the Controllers}
\label{subsec:controller-comparison}

The three controllers have been compared in simulation while tracking
a circle, a square, and a figure eight. The outcomes:
\begin{itemize}
    \item \textbf{Approximate linearization}: only \textbf{local}
          stability, and only for constant \( v_d, \omega_d \). It
          converges to zero error on a circle; on a square it produces
          a new transient at each corner; on a figure eight local
          stability is not guaranteed, though the observed performance
          is good.
    \item \textbf{Static input-output linearization}: the steady-state
          error norm equals \( b \), because it is the point \( B \)
          that converges to the desired trajectory, while the contact
          point stays at distance \( b \) from it. The displaced output
          provides a \textbf{lookahead} behavior when tracking a
          square. Reducing \( b \) reduces the steady-state error, but
          the steering velocity increases, since the term \( 1/b \)
          appears in \( T^{-1} \).
    \item \textbf{Dynamic input-output linearization}: zero
          steady-state error, reasonable velocities, and faster
          transients (for example at the corners of the square).
\end{itemize}

\section{Takeaways}
\label{sec:motion-control-takeaways}

\begin{enumerate}
    \item Motion control for wheeled mobile robots is designed at the
          \textbf{kinematic level}: the dynamic model is canceled by
          the low-level velocity loops.
    \item In \textbf{trajectory tracking} the desired Cartesian
          trajectory is mapped to the desired state and inputs via
          flatness, and a feedback action is added on the state or
          output error.
    \item \textbf{Approximate linearization} of the tracking error
          gives a linear time-varying error dynamics. Local asymptotic
          stability follows from the indirect Lyapunov method, but
          only for constant reference inputs, and the controller
          requires \( v_d \neq 0 \).
    \item \textbf{Exact input-output linearization} is not possible
          for the Cartesian position of the unicycle: the decoupling
          matrix is singular. Two fixes: displace the output along
          the sagittal axis (static feedback, with zero dynamics and
          steady-state error proportional to the displacement), or
          extend the dynamics by treating \( v \) as a state (dynamic
          feedback, no zero dynamics, requires \( v \neq 0 \)).
    \item All controllers require a \textbf{persistent} trajectory:
          the nonholonomic singularity at \( v = 0 \) is structural
          and cannot be removed by feedback.
\end{enumerate}